\noindent 
\textbf{Abstract -- } The emergence of Large Language Models (LLMs) has popularized ``vibe coding,'' a paradigm where software is developed primarily through natural language prompting. While this approach lowers entry barriers and accelerates initial code generation, it frequently precipitates a ``Debugging Paradox'' for novice programmers, who find themselves unable to troubleshoot AI-generated codebases they did not cognitively construct. This study empirically investigates the impact of vibe coding on the cognitive integrity and problem-solving resilience of undergraduate students through a controlled observational experiment. Participants performed an API integration task within a development environment featuring a deliberate, undocumented dependency mismatch designed to trigger a failure in AI-generated solutions. Results demonstrate a pronounced ``cognitive erosion,'' where students bypass deep learning in favor of passive AI delegation. This behavior led to a ``debugging trap'' characterized by excessive ``prompt-looping'' and a pervasive ``illusion of competence'' driven by the AI's rapid output. Furthermore, the findings highlight ``architectural blindness,'' as participants struggled to comprehend the systemic dependencies within their generated code. The study concludes that over-reliance on vibe coding may fundamentally hinder the development of mental models, suggesting that foundational manual programming skills remain critical for navigating the limitations of AI-assisted development.

\smallskip
\smallskip

\textbf{Keywords -- Vibe Coding; Large Language Models; Debugging Paradox; Programming Education; Cognitive Offloading; AI-Assisted Development}